\chapter{Tools: The Execution Layer}
\label{ch:tools}

\epigraph{``You're gonna need a bigger boat.''}{Chief Brody, \emph{Jaws} (1975)}

Brody says it after seeing the thing he has been theorising about, at actual
scale, for the first time. Everybody has this moment with tool catalogues. You
design six tools, they work beautifully, you ship. Eighteen months later there
are forty-one, three of them do nearly the same thing, and the model has started
picking the wrong one about a fifth of the time and nobody can say when that
started.

This chapter is about not getting there. Tools are the primitive that does the
work, and they are the primitive where design mistakes compound fastest, because
their cost is charged on every request forever.

\section{Mechanics first}

Two methods. That is the entire surface.

\subsection{\texttt{tools/list}}

\begin{wire}
{
  "jsonrpc": "2.0",
  "id": 1,
  "result": {
    "resultType": "complete",
    "tools": [
      {
        "name": "assess_account_risk",
        "title": "Assess account risk",
        "description": "Score one commercial account for credit risk...",
        "inputSchema": { "type": "object", "properties": { ... } },
        "outputSchema": { "type": "object", "properties": { ... } },
        "annotations": { "readOnlyHint": true, "idempotentHint": true }
      }
    ],
    "nextCursor": "50",
    "ttlMs": 300000,
    "cacheScope": "public"
  }
}
\end{wire}

Three fields at the bottom that did not exist in 2025 and matter a great deal.
\texttt{nextCursor} means catalogues paginate. \texttt{ttlMs} and
\texttt{cacheScope} mean the client can stop asking. Chapter~\ref{ch:resources}
covers the caching model; here the relevant fact is that a well-behaved client fetches your catalogue once per TTL, not once per task.

There is also a requirement that reads like a footnote and is not:

\keyidea{Return tools in a deterministic order. The same order, every time, when
the underlying set has not changed. Unstable ordering silently destroys the
provider's prompt cache, and nothing in your telemetry will tell you.}

That last part is what makes it nasty. Prompt caching keys on a stable prefix.
If your catalogue serialises in a different order on each request, the prefix
changes, every cache lookup misses, and you pay full price for prompt processing
on every turn. Your server looks fine. Your latency is quietly 30\,\% worse.
Meridian sorts by name and asserts it:

\begin{py}
def test_tools_list_order_is_stable(self):
    """Unstable ordering silently destroys the provider's prompt cache."""
    server = tiny_server()
    for i in range(6):
        server.add_tool(_noop_tool(f"tool_{i}"))
    orders = [
        [t["name"] for t in server.handle(
            request("tools/list", req_id=i))["result"]["tools"]]
        for i in range(5)
    ]
    self.assertEqual(len(set(map(tuple, orders))), 1)
\end{py}

If you build your catalogue by iterating a Python dictionary, a Go map, or a
database query with no \texttt{ORDER BY}, you are one refactor away from this
bug. Sort explicitly.

\subsection{\texttt{tools/call}}

\begin{wire}
{
  "jsonrpc": "2.0",
  "id": 2,
  "method": "tools/call",
  "params": {
    "name": "assess_account_risk",
    "arguments": { "accountId": "ACC-1042", "exposureUsd": 250000 }
  }
}
\end{wire}

And back:

\begin{wire}
{
  "jsonrpc": "2.0",
  "id": 2,
  "result": {
    "resultType": "complete",
    "content": [{ "type": "text", "text": "{\"score\":55.4,...}" }],
    "structuredContent": { "accountId": "ACC-1042", "score": 55.4,
                           "band": "elevated", "drivers": [...] },
    "isError": false
  }
}
\end{wire}

The result carries the same data twice, deliberately: \texttt{structuredContent} for anything that will parse it,
\texttt{content} for backwards compatibility and for models that do better with
text. Chapter~\ref{ch:apps} shows the third consumer, an MCP App, which reads
\texttt{structuredContent} and never puts it in the model's context at all.

\section{Designing tools a model can actually use}

Now the part that matters more than the mechanics.

\subsection{Stop mirroring your REST API}

The single most common failure. You have an API. It has endpoints. Wrapping each
endpoint as a tool is mechanical, fast, and completable in an afternoon, which is
exactly why so many people do it.

Meridian shipped this way. You can still run it:

\begin{sh}
python3 -m meridian.serve risk --fat
\end{sh}

It exposes \texttt{get\_account} and \texttt{list\_account}, then
\texttt{create\_exposure}, then \texttt{update\_covenant}, then
\texttt{delete\_collateral}, and two dozen more in the same shape. Every one is
a faithful wrapper around a REST endpoint. Every one works.

Here is what it costs.

\bookfig[0.98]{catalogue-cost}{The same server, two catalogues. Left: what the
catalogue costs. Right: what it costs by turn eight, which is the number that
should change your mind.}{catalogue-cost}

\begin{measurebox}{REST mirror versus task-shaped tools}
\begin{tabular}{@{}lrr@{}}
\toprule
\thead{} & \thead{REST mirror} & \thead{Task-shaped} \\
\midrule
Tools in catalogue          & 38    & 10 \\
Tokens per model turn       & 5{,}617 & 1{,}246 \\
Extra tokens per turn       & +4{,}371 & \\
Extra tokens per task       & +8{,}738 & \\
Wasted input cost per 1,000 tasks & \$26.21 & \\
\bottomrule
\end{tabular}

\vspace{6pt}
\footnotesize
Meridian, four servers, measured by \texttt{make bench}. The task-shaped
catalogue does strictly more useful work.
\end{measurebox}

Twenty-six dollars per thousand tasks, for nothing. And that is only the money.
The real damage is upstream: the model now chooses among 38 similar options
instead of 10 distinct ones, and selection accuracy falls, and a wrong choice
costs a whole extra loop iteration to detect and recover from.

Why does the REST shape fail? Because REST is organised around \emph{resources
and verbs}, and a planner reasons in \emph{tasks and outcomes}. To answer ``is
this account risky'', a model given a REST mirror must plan: get the account,
list its exposures, list its covenants, fetch the rating, then combine. Five
round trips, five model turns, five chances to go wrong.

Given a task-shaped tool it calls \texttt{assess\_account\_risk} once.

\keyidea{Design tools for the plan, not for the database. The right question is
not ``what are my entities'', it is ``what does the user want done, and what is
the smallest number of calls that does it.''}

\subsection{The three rewrites that fix most catalogues}

\textbf{Collapse read-then-combine chains.} Any sequence the model performs
identically every time is a tool you have not written yet. If your traces show
\texttt{get\_x} always followed by \texttt{get\_y} then \texttt{compute\_z},
ship \texttt{compute\_z\_for\_x}.

\textbf{Add a batch variant of anything called in a loop.} Meridian has
\texttt{assess\_account\_risk} for one account and
\texttt{rank\_portfolio\_risk} for many. Without the second, a portfolio review
of forty accounts is forty round trips and forty model turns. With it, one.
The description points at it explicitly, because the model reads descriptions:

\begin{py}
# sketch: elided from the shipped handler
@server.tool(
    "rank_portfolio_risk",
    "Rank a set of accounts by risk score in one call. Prefer this over "
    "calling assess_account_risk repeatedly.",
    ...
)
\end{py}

\textbf{Delete tools nobody calls.} You have telemetry. Any tool with no calls
in ninety days is costing you tokens on every request in exchange for nothing.
Delete it. If somebody objects, the objection is usually about the endpoint
existing, not about the tool existing, and those are different things.

\subsection{Descriptions are prompt engineering}

The description is not documentation. It is the only thing the model sees when
deciding, and it is charged per turn. So it should be short and it should
contain the information that distinguishes this tool from its neighbours.

Bad:

\begin{plain}
Get a account record via the core banking API. Wraps GET /v2/accounts.
Accepts an optional id, an optional body, and standard pagination
parameters. Returns the raw account representation as stored by the
system of record, including audit fields, soft-delete markers, and
denormalised references to related entities.
\end{plain}

Fifty-one tokens describing the transport and the storage layer, neither of
which is a fact the model can act on.

Good:

\begin{plain}
Score one commercial account for credit risk. Returns a 1-99 score, a
band, and the weighted factors behind it.
\end{plain}

Twenty-three tokens, all of them about what the model gets back.

Three rules that survive contact with reality:

\begin{enumerate}[leftmargin=1.6em, itemsep=3pt]
\item \textbf{Lead with the outcome}, not the mechanism. The model is choosing
between outcomes.
\item \textbf{Name the discriminator} when tools are adjacent. ``Prefer this
over calling X repeatedly'' is worth its tokens.
\item \textbf{Do not document your transport.} The model cannot use
``wraps GET /v2/accounts'' and you pay for it forever.
\end{enumerate}

\section{Schemas}

JSON Schema 2020-12 by default, and \texttt{2026-07-28} loosened the restriction
so you can use the full dialect. That freedom comes with two obligations.

\subsection{The \texttt{\$ref} rule, which is a security rule}

\begin{dangerbox}{Never dereference a remote \texttt{\$ref}}
JSON Schema permits \texttt{\$ref} to an absolute URI. A validator that fetches
it is a server-side request forgery primitive: a hostile tool definition points
\texttt{\$ref} at \texttt{http://169.254.169.254/} and your validator obligingly
retrieves cloud instance credentials.

The specification makes this a MUST NOT by default. An opt-in mode may exist,
but it must be off by default, and it should enforce a host allowlist, reject
loopback and private addresses, apply timeouts and size limits, and log every
URI it touched.
\end{dangerbox}

\begin{py}
if "$ref" in schema and str(schema["$ref"]).startswith(("http://", "https://")):
    raise errors.InvalidParams(f"{where}: remote $ref is not dereferenced")
\end{py}

A schema that fails to validate because of an unresolved external \texttt{\$ref}
must be \emph{rejected}, not silently treated as permissive. Failing open here
means an attacker can disable your validation by pointing at a URL you will not
fetch.

\subsection{The composition rule, which is a denial-of-service rule}

\texttt{anyOf}, \texttt{oneOf}, \texttt{allOf}, \texttt{if}/\texttt{then}, and
\texttt{\$defs} are expressive and expensive. Nested deeply enough, validation
cost explodes combinatorially, and a hostile schema becomes a CPU bomb against
whoever validates it.

So bound it: a maximum depth, a cap on subschemas, or a per-validation time
budget. Meridian caps depth at twelve, which no honest tool schema approaches:

\begin{py}
if depth > 12:
    raise errors.InvalidParams(f"{where}: schema nesting exceeds the depth limit")
\end{py}

\subsection{Output schemas earn their keep}

Optional, and you should use them anyway.

Without one, a model receives text and has to parse it. Sometimes it parses
wrong, which costs a retry, which costs a model turn. With one, the client
validates \texttt{structuredContent} before it ever reaches the model, and a malformed response is caught at the boundary, before it can confuse the
planner two turns later.

\begin{py}
output_schema={
    "type": "object",
    "properties": {
        "accountId": {"type": "string"},
        "score": {"type": "number"},
        "band": {"type": "string",
                 "enum": ["low", "moderate", "elevated", "high"]},
        "drivers": {"type": "array", "items": {...}},
    },
    "required": ["accountId", "score", "band", "drivers"],
}
\end{py}

That \texttt{enum} on \texttt{band} is doing real work. It tells the model the
complete set of possible answers, which means it will not invent
\texttt{"very high"} and it will not write branching logic for a fifth case that
cannot occur.

\subsection{Annotations, and why they are untrusted}

\begin{py}
annotations={"readOnlyHint": True, "idempotentHint": True}
\end{py}

Useful hints. A host can auto-approve read-only tools and prompt for the rest,
which is exactly what Meridian's default consent policy does.

But the specification is blunt: clients must treat annotations as untrusted
unless the server is trusted. A compromised server marking
\texttt{delete\_everything} as \texttt{readOnlyHint: true} would sail through a
host that believes annotations.

\begin{py}
def check(self, qualified_name: str, tool: dict, arguments: dict) -> str:
    if qualified_name in self.denylist:
        decision = ConsentDecision.DENY
    elif self.auto_allow_read_only and \
            (tool.get("annotations") or {}).get("readOnlyHint"):
        decision = ConsentDecision.ALLOW
    ...
\end{py}

Meridian does auto-allow on \texttt{readOnlyHint}, and that is a defensible
default \emph{for servers you trust}. Chapter~\ref{ch:threats} shows what
happens when you extend the same courtesy to a server you should not.

\section{Errors: the distinction that decides whether a model recovers}

This is the most consequential paragraph in the chapter, so here it is plainly.

There are two error mechanisms, they are not interchangeable, and mislabelling
one as the other either burns tokens or kills tasks.

\textbf{Protocol errors} are JSON-RPC errors. They mean the request was
structurally wrong: unknown tool, malformed parameters, server on fire. Models
are unlikely to recover from these, because the problem is not in the arguments,
it is in the shape.

\begin{wire}
{ "jsonrpc": "2.0", "id": 3,
  "error": { "code": -32602, "message": "Unknown tool: invalid_tool_name" } }
\end{wire}

\textbf{Tool execution errors} are \emph{successful} results with
\texttt{isError: true}. They mean the call was well-formed and the work failed
for a reason the model can understand and act on.

\begin{wire}
{ "jsonrpc": "2.0", "id": 4,
  "result": {
    "resultType": "complete",
    "content": [{ "type": "text",
      "text": "Invalid departure date: must be in the future. Today is 2026-08-08." }],
    "isError": true } }
\end{wire}

Read that message. It says what was wrong, and it gives the model what it needs
to fix it. That is the whole design goal.

\begin{table}[htbp]
\centering
\small
\begin{tabularx}{\textwidth}{@{}XllX@{}}
\toprule
\thead{Situation} & \thead{Mechanism} & \thead{Code} & \thead{Because} \\
\midrule
Unknown tool name          & Protocol & \texttt{-32602} & The catalogue is wrong, not the arguments \\
Missing required argument  & Protocol & \texttt{-32602} & Schema violation; caught before your code runs \\
Account id does not exist  & Tool     & \texttt{isError} & The model can try a different id \\
Date in the past           & Tool     & \texttt{isError} & The model can pick another date \\
Downstream API timed out   & Tool     & \texttt{isError} & The model can retry or route around \\
Insufficient permissions   & Tool     & \texttt{isError} & The model can explain, or take another path \\
Server crashed             & Protocol & \texttt{-32603} & Nothing the model does will help \\
\bottomrule
\end{tabularx}
\caption{Which mechanism, and why. The rule of thumb: if a smart colleague
could act on the message, it is a tool execution error.}
\label{tab:error-choice}
\end{table}

Meridian's version, and note how much work the message does:

\begin{py}
account = ACCOUNTS.get(account_id)
if account is None:
    # A tool execution error, not a protocol error. The model can read
    # this, notice it used a stale id, and try a different one.
    return tool_error(
        f"No account {account_id}. Account ids look like ACC-1000 "
        f"through ACC-{1000 + len(ACCOUNTS) - 1}."
    )
\end{py}

It states the failure and the valid range. A model reading that fixes its next
call. A model reading ``Error: not found'' does not, and asks the user, and now
a recoverable situation has become a conversation.

\begin{py}
def test_bad_business_input_is_a_tool_execution_error(self):
    resp = build_server().handle(request(
        "tools/call",
        {"name": "assess_account_risk", "arguments": {"accountId": "ACC-9999"}}))
    self.assertNotIn("error", resp)
    self.assertTrue(resp["result"]["isError"])
    # And it must be actionable, not just "invalid".
    self.assertIn("ACC-1000", resp["result"]["content"][0]["text"])
\end{py}

\section{The interception loop}

Between the model's intent and the server's execution sits the host, and the
places it can intervene are the places your security lives.

\bookfig[0.98]{interception-loop}{Every dotted line is a hook, and every hook is
a place to say no. A host that forwards intents untouched has not implemented
MCP.}{interception-loop}

Four outbound hooks and one inbound one:

\begin{description}[leftmargin=1.4em, style=nextline, itemsep=3pt]
\item[Schema validation] Reject malformed arguments before they reach a server.
Cheap, and it turns a class of server bug into a client-side error.
\item[Consent] Show the user the arguments, not just the tool name. ``Call
\texttt{send\_email}?'' is not consent. ``Send email to \texttt{whom}, saying
\texttt{what}?'' is.
\item[Audit] Who, what, when, with which arguments, and what came back. Designed
for a regulator, not for your debugger. Chapter~\ref{ch:access-control} covers
what that means concretely.
\item[Authorization] Scope checks and row-level policy, on the server side,
where the data is.
\item[Result validation] The return path. Validate against
\texttt{outputSchema}, then treat the content as data. Never as instructions.
\end{description}

That last one is the one people skip, and it is the one Chapter~\ref{ch:threats}
is about.

\section{Parallel fan-out}

When a model emits three tool calls in one turn, it is telling you those three
are independent. Honour that.

\begin{py}
def call_tools_parallel(self, calls: list[dict]) -> list[dict]:
    """Fan out independent calls.

    The model asked for these together, which is its way of saying none of
    them depends on the others. Running them in sequence throws that
    information away and pays the sum of the latencies instead of the max.
    """
    if len(calls) <= 1:
        return [self.call_tool(c["name"], c.get("arguments")) for c in calls]

    with concurrent.futures.ThreadPoolExecutor(max_workers=len(calls)) as pool:
        futures = [pool.submit(self.call_tool, c["name"], c.get("arguments"))
                   for c in calls]
        ...
\end{py}

Chapter~\ref{ch:primer} measured this: a consistent 2.8$\times$ on three calls,
holding steady as per-call latency rises from 2\,ms to 40\,ms. The ceiling is
3$\times$, and the gap is thread scheduling plus the slowest of the three
servers.

Two things the host must get right for this to work at all.

\textbf{Preserve ordering.} Results must come back in the order the calls were
issued, whatever order they complete in. Meridian iterates futures in submission order for exactly this reason.

\textbf{Never fail the batch on one failure.} One tool erroring should produce
one error result, not an exception that discards the other two successful calls
you already paid for.

\begin{notebox}{You cannot fan out what the model serialised}
If your tool design forces \texttt{get\_x} then \texttt{get\_y} then
\texttt{combine}, the model \emph{cannot} parallelise, because each step depends
on the last. Parallel fan-out is available only to catalogues whose tools are
independent, which is another way of saying tool design and latency are the same
problem.
\end{notebox}

\section{Meridian's catalogue, and why it is ten tools}

\begin{table}[htbp]
\centering
\small
\begin{tabularx}{\textwidth}{@{}llX@{}}
\toprule
\thead{Server} & \thead{Tool} & \thead{Why it exists} \\
\midrule
risk & \texttt{assess\_account\_risk} & The single-account question, one call \\
risk & \texttt{rank\_portfolio\_risk} & So a 40-account review is 1 trip, not 40 \\
risk & \texttt{underwrite\_loan}      & Long-running, returns a task (\chapref{tasks}) \\
compliance & \texttt{review\_transaction} & Needs a named officer (\chapref{mrtr}) \\
compliance & \texttt{search\_guidance}   & Reference text. Untrusted (\chapref{threats}) \\
compliance & \texttt{export\_audit\_log} & For the regulator, not the debugger \\
fraud & \texttt{screen\_account}   & Fast, read-only, safe to run in parallel \\
fraud & \texttt{explain\_signal}   & So the model can explain a signal \\
marketdata & \texttt{get\_reference\_curve} & Public, cacheable, called constantly \\
marketdata & \texttt{price\_facility}      & The combine step, done server-side \\
\bottomrule
\end{tabularx}
\caption{Ten tools across four servers. Each is a task somebody wants done, and
none is an entity with CRUD hung off it.}
\label{tab:meridian-tools}
\end{table}

The \texttt{price\_facility} row deserves a note, because it is the pattern
generalised. Pricing needs a risk score, a tenor, and a sector spread. You
\emph{could} expose \texttt{get\_curve}, \texttt{get\_sector\_spread}, and
\texttt{get\_credit\_spread} and let the model do arithmetic. Three round trips,
three model turns, and a model doing floating-point maths in its head.

Instead the server does the combining and returns the answer. One call. And the
arithmetic is done by a computer, which is generally better at it.

\section{What to remember}

\begin{itemize}[leftmargin=1.6em, itemsep=3pt]
\item Two methods total. The design work is all in what you put in them.
\item Deterministic ordering, or you silently lose the provider's prompt cache.
\item Task-shaped, not REST-shaped. Meridian's REST mirror cost 4,371 extra
tokens per turn and \$26 per thousand tasks in pure waste.
\item Descriptions lead with the outcome, name the discriminator, and never
mention your transport.
\item Never dereference remote \texttt{\$ref}s. Bound schema depth.
\item Output schemas eliminate parse-retry loops. Use \texttt{enum} generously.
\item Protocol errors are for malformed requests. Tool execution errors are for
things the model can fix, and the message should tell it how.
\item Fan out what the model batched, keep the ordering, and do not let one
failure discard the batch.
\end{itemize}

Next: resources, which is where the caching model lives and where the token
budget is most often wasted.
